\subsection{高维情形}\label{sec:20b}
假设 $n\ge2$。此时，可将\stref{引理}{1.9}与适当的 Hochschild 上同调群消失结果结合，证明任意上同调为 $A_m^n$ 的 $\Ai$ 范畴必为形式的，即拟同构于 $A_m^n$ 上的平凡 $\Ai$ 结构，后者的所有阶数大于 $2$ 的复合映射都为零。这已在\bcite{134}中完成。下面采用更具体的方法证明这一结果，避免使用抽象分类理论。在几何方面，这依赖于第\sr{section:18}{18 节}的材料，因此需要技术假设 $\operatorname{char}(\K)\ne2$。

取函数 $\pi:X\to\CC$，$\pi(x)=x_0$。再次经过技术调整后，可将它视为 Lefschetz 纤维化，将 $(\Psi,\pi)$ 视为双纤维化。此外，可显式分析 $\pi$ 的纤维与平行移动映射\bcite{85}：每个光滑纤维都辛同构于余切丛 $T^*S^{n-1}$；每条消失路径都给出同一个消失圈，即零截面
\SourcePage{314}
$S^{n-1}\subset T^*S^{n-1}$。因此，每条连接两个临界值的路径自然都是匹配路径。特别地，$V_j$ 是一链匹配路径 $\mu_j$ 对应的匹配圈，见\sr{fig:20.1}{图~20.1}。
\SourceFigure{figure-20-1}{0.350999622}{20.1}

此后固定 $m$ 及 $n\ge2$，写 $\F=\F_m^n$、$A=A_m^n$。若只考虑余切丛中的零截面 $S^{n-1}$，则相应 Fukaya $\Ai$ 代数拟同构于 $\K\cdot e\oplus\K\cdot f$，其中
\begin{equation}\label{eq:20.4}
\begin{gathered}
|e|=0,\qquad |f|=n-1,\\
\mu^2(e,e)=e,\qquad\mu^2(e,f)=(-1)^{n-1}f,\qquad\mu^2(f,e)=f,
\end{gathered}
\end{equation}
其余所有乘积都为零。\znote{306}{(20.4) 中 $\mu^2(e,f)=(-1)^{n-1}f$ 而 $\mu^2(f,e)=f$，并非笔误：它等价于 $\mu^2(e,f)=(-1)^{|f|}\mu^2(f,e)$，即单位元 $e$ 在左侧作用时带符号 $(-1)^{|x|}$。这是本书 $\Ai$ 符号约定（输入按圆盘边界顺序排列）的产物，不能按严格含幺代数的 $e\cdot f=f$ 去“改正”，否则\eref{20.5}的复合及后文各式的符号会整体出错。另一种看法：$e$ 是单位、$f$ 是 $S^{n-1}$ 的基本类，$\mu^2(f,e)=f$ 即杯积，而 $\mu^2(e,f)$ 上的符号正是把 $e$ 从右端移到左端所付的代价。[来源：A，原书 (20.4) 及其后的 (20.5)。]}在上同调层次上，这是显然的，因为考察的只是
\[
HF^*(S^{n-1},S^{n-1})=H^*(S^{n-1};\K).
\]
在上链层次上，所需消失结果来自次数考察，至少在把 $\Ai$ 结构做成严格含幺之后是如此。另一种方法是使用\bcite{58}的结果，以及 $S^{n-1}$ 在传统代数拓扑意义下为形式的这一事实。下一步是写出 $\pi$ 的一个消失圈基所对应的有向 $\Ai$ 范畴 $\B^\to$；该基当然只由 $m+1$ 个 $S^{n-1}$ 的副本构成。这很直接：将对象记为 $(Z_0,\ldots,Z_m)$，则态射为
\begin{equation}\label{eq:20.5}
\hom_{\B^\to}(Z_i,Z_j)=\begin{cases}
\K\cdot e_{ji}\oplus\K\cdot f_{ji},&i<j,\\
\K\cdot e_{ii},&i=j,\\
0,&i>j,
\end{cases}
\end{equation}
次数与复合均同\eref{20.4}。取 $\Tw\B^\to$ 的对象 $(C_1,\ldots,C_m)$，定义为
\begin{equation}\label{eq:20.6}
C_j=\Cone(e_{j,j-1}).
\end{equation}
\SourcePage{315}
\stref{命题}{18.21}说明，由这些对象构成的全 $\Ai$ 子范畴拟同构于 $\F$。容易验证，在上同调层次上，这恢复了 $A$。更重要的是，该过程给出显式的链层次模型。最后，这个模型还具有一个有用的额外对称性，它并不先验地来自几何。

为简化记号，使用\stref{笔记}{3.26}中的扭曲复形版本。给 $\B^\to$ 再赋予长度分次，使 $e_{ji}$ 的长度为零，$f_{ji}$ 的长度为二；这一名称稍后会得到解释。于是可引入\emph{齐次加性完备化} $\Sigma^{\mathrm{ho}}\B^\to$，其对象为有限和
\begin{equation}\label{eq:20.7}
\bigoplus_i S^{-\sigma_i}\Lambda^{-\tau_i}Y_i,
\end{equation}
其中 $\Lambda$ 为使长度分次降低一的形式运算。任意两个这样的对象之间的态射，便带有按次数与长度给出的双分次。同样，有齐次扭曲复形的 $\Ai$ 范畴 $\Tw^{\mathrm{ho}}\B^\to$，其对象是带保持长度分次的微分的表达式\eref{20.7}。\znote{307}{为什么要多引入一个长度分次：次数分次无法表达“$\Cone(e_{j,j-1})$ 应记作长度为 $j$ 的对象”这一信息。把 $e_{ji}$ 的长度定为 $0$、$f_{ji}$ 定为 $2$，再用形式运算 $\Lambda$ 降低长度，就能把 (20.6) 写成 $S\Lambda^jZ_{j-1}\oplus\Lambda^jZ_j$，使“关于长度分次齐次”的 $\Ai$ 结构有意义；这样用\stref{命题}{1.12}消去 $\mu^1$ 时才不会破坏双分次，最后才能断言 $d>2$ 的高阶 $\mu^d_{\A}$ 必为零。[来源：A，原书本节。]}例如，可将\eref{20.6}变成如下齐次扭曲复形：
\[
C_j=\Lambda^j\Cone(e_{j,j-1})=S\Lambda^jZ_{j-1}\oplus\Lambda^jZ_j,\qquad
\delta_{C_j}=\begin{pmatrix}0&0\\e_{j,j-1}&0\end{pmatrix}.
\]
如前，取 $\Tw^{\mathrm{ho}}\B^\to$ 中由这些对象构成的全 $\Ai$ 子范畴，它是双分次 $\Ai$ 范畴。在上同调层次上，这一额外分次确实恢复了 $A$ 作为\eref{20.2}的路径代数的商所继承的长度分次。

对前面构造的全 $\Ai$ 子范畴应用\stref{命题}{1.12}，以消去 $\mu^1$。可用保持双分次的方式做到这一点，结果是 $A$ 上关于长度分次齐次的 $\Ai$ 结构 $\A$。考察可知，此时高阶复合映射 $\mu_{\A}^d$，$d>2$，必为零。另一方面，由构造，这一 $\Ai$ 结构拟同构于 $\F$，所以 $\F$ 为形式的。
